\setupcolors[state=start]
\setupcolor[hex]
\definecolor[darkblue][h=003366]
\definecolor[darkgray][h=e6e6e6]
\definecolor[redline][h=d9534f]
\definecolor[blueline][h=5bc0de]
\definecolor[greenline][h=5cb85c]

\setuphead[section][style=\bfa, color=darkblue, before=\blank, after=\blank]
\setuphead[subsection][style=\bf, color=black, before={\blank[2*big]}, after=\blank]
\setuphead[subsubsection][style=\bf, color=darkblue, before=\blank, after=\blank]
\setuppapersize[A4][A4]
\setupbodyfont[11pt]
\setuplayout[width=middle, height=middle, topspace=2cm, bottomspace=2cm]
\mainlanguage[es]

\starttext

\title{\bf Modelos OPP v5 y heurísticas HC-4 y ACO-4:\crlf definiciones y comparación}
\blank[big]

%%%%%%%%%%%%%%%%%%%%%%%%%%%%%%%%%%%%%%%%%%%%%%%%%%%%%%%%%%%%%%%%%%%%
\startsection[title={Introducción}]

Este documento define los dos modelos exactos de v5 y las dos heurísticas que resuelven el
mismo problema, y los compara con los resultados de los experimentos. Los cuatro métodos
resuelven el mismo problema de {\em order-picking}: dado un conjunto de lotes
de pedidos, un grafo de locaciones de almacén y un conjunto de pickers heterogéneos en
velocidad, determinar simultáneamente la ruta de recolección de cada lote, la asignación de
lotes a pickers y la secuenciación temporal, minimizando una función objetivo que combina
tiempos de completitud con penalizaciones de tardanza y adelanto, sujeto a una restricción
de cadena de frío.

Las dos formulaciones exactas difieren en la estrategia de resolución:

\startitemize
\item {\bf OPP v5.1 (monolítico)}: un único MIP que decide ruta, asignación y secuenciación
en conjunto. El tiempo de ruta $r_b$ es una {\em variable} de decisión.
\item {\bf OPP v5.2 (dos etapas)}: la Etapa~1 resuelve el ruteo óptimo de cada lote de
forma independiente (minimizando el tiempo de ruta); la Etapa~2 recibe el tiempo de ruta
como {\em dato} fijo y resuelve asignación, secuenciación y cadena de frío.
\stopitemize

Las dos formulaciones tienen el mismo valor óptimo: la ruta más corta domina cualquier otra porque el
objetivo no contiene un término de tiempo de ruta y una ruta más corta solo relaja las
restricciones de duración y cadena de frío. Sin embargo, el conjunto de soluciones óptimas
difiere: el modelo monolítico puede devolver rutas no mínimas que alcanzan el mismo
objetivo.

Las dos heurísticas (sección \in[sec:heuristicas]) buscan sobre el mismo espacio de
soluciones (una ruta por lote, una asignación de lotes a pickers y una secuencia por picker)
con el mismo objetivo, pero sin garantía de optimalidad:
\startitemize
\item {\bf OPP\_hc\_4 (HC-4)}: búsqueda local estocástica ({\em Hill Climbing}) que parte de
rutas construidas por varios métodos y mejora rutas, asignación y secuencia con movimientos
aleatorios.
\item {\bf OPP\_aco\_4 (ACO-4)}: colonia de hormigas ({\em Ant Colony Optimization}) que
construye rutas y asignaciones con feromonas y aplica búsqueda local a la mejor solución.
\stopitemize

La sección \in[sec:comparacion] compara los cuatro métodos: qué instancias son comparables,
cuánto se alejan las heurísticas del óptimo, cuánto tardan, y qué obtiene el modelo exacto si
recibe el mismo tiempo que usó la heurística.

Los archivos fuente son:
\startitemize
\item \type{models/opp_v5_1_mono/opp_model_v5_1_mono.zpl}
\item \type{models/opp_v5_2_split/opp_model_v5_2_split_stage1_routing.zpl}
\item \type{models/opp_v5_2_split/opp_model_v5_2_split_stage2_scheduling.zpl}
\item \type{models/opp_hc_4/opp_hc_4.py}
\item \type{models/opp_aco_4/opp_aco_4.py}
\stopitemize

\stopsection


%%%%%%%%%%%%%%%%%%%%%%%%%%%%%%%%%%%%%%%%%%%%%%%%%%%%%%%%%%%%%%%%%%%%
\startsection[title={Notación común}]

Los dos modelos exactos usan los mismos datos. En esta sección se definen una sola vez.

\startsubsection[title={Conjuntos}]
\startformula
\startalign[n=2, align={right,left}]
    \NC \mathcal{N} \NC \quad \text{órdenes de clientes} \NR
    \NC \mathcal{B} \NC \quad \text{lotes; cada orden pertenece a un único lote} \NR
    \NC \mathcal{P} \NC \quad \text{pickers} \NR
    \NC \mathcal{V} \NC \quad \text{locaciones del almacén, incluidos el depósito } v_{start} \text{ y el nodo final } v_{end} \NR
    \NC \mathcal{E} \NC \quad \text{arcos del grafo de locaciones: pares } (i,j) \text{ por los que se puede ir de } i \text{ a } j \NR
\stopalign
\stopformula

Para escribir las restricciones de forma compacta se usan además:
\startformula
\startalign[n=2, align={right,left}]
    \NC \mathcal{V}^{int} \NC = \mathcal{V} \setminus \{v_{start}, v_{end}\} \quad \text{(locaciones intermedias)} \NR
    \NC \delta^{+}(i) \NC = \{(i,j) \in \mathcal{E}\} \quad \text{(arcos que salen de } i\text{)} \NR
    \NC \delta^{-}(i) \NC = \{(j,i) \in \mathcal{E}\} \quad \text{(arcos que entran a } i\text{)} \NR
    \NC \mathcal{N}_b \NC = \{ n \in \mathcal{N} : B_{nb} = 1 \} \quad \text{(órdenes del lote } b\text{)} \NR
\stopalign
\stopformula
\stopsubsection

\startsubsection[title={Parámetros}]

{\bf Datos de órdenes, lotes y almacén}
\startformula
\startalign[n=2, align={right,left}]
    \NC A_{ni} \in \{0,1\} \NC \quad 1 \text{ si la orden } n \text{ tiene productos en la locación } i \NR
    \NC B_{nb} \in \{0,1\} \NC \quad 1 \text{ si la orden } n \text{ pertenece al lote } b \NR
    \NC C_n \NC \quad \text{capacidad (cantidad a recoger) de la orden } n \NR
    \NC T^{due}_n \NC \quad \text{fecha de entrega de la orden } n \NR
    \NC D_{ij} \NC \quad \text{distancia del arco } (i,j) \in \mathcal{E} \NR
    \NC S_i \NC \quad \text{tiempo de servicio en la locación } i \NR
\stopalign
\stopformula

{\bf Tiempos y pickers}
\startformula
\startalign[n=2, align={right,left}]
    \NC T^{travel} \NC \quad \text{tiempo por unidad de distancia, a velocidad de referencia} \NR
    \NC T^{setup} \NC \quad \text{tiempo de preparación de cada lote} \NR
    \NC T^{pick} \NC \quad \text{tiempo de picking por unidad de capacidad} \NR
    \NC V_p \NC \quad \text{factor de velocidad del picker } p \text{ (} 1{,}0 = \text{referencia; } 2{,}0 = \text{camina el doble de rápido)} \NR
    \NC T^{cold} \NC \quad \text{tiempo máximo que un lote puede estar fuera de la cámara de frío} \NR
    \NC \varepsilon \NC \quad \text{tolerancia numérica del chequeo de frío } (10^{-9}) \NR
\stopalign
\stopformula

{\bf Pesos del objetivo}
\startformula
\startalign[n=2, align={right,left}]
    \NC W^{c} \NC \quad \text{peso de los tiempos de completitud de los lotes} \NR
    \NC W^{pen} \NC \quad \text{peso del total de penalizaciones} \NR
    \NC W^{e},\; W^{t} \NC \quad \text{peso del adelanto y de la tardanza dentro de las penalizaciones} \NR
\stopalign
\stopformula

{\bf Constantes grandes (Big-M)}
\startformula
\startalign[n=2, align={right,left}]
    \NC M \NC \quad \text{constante grande general, dato del escenario} \NR
    \NC M^{seq} \NC \quad \text{constante grande de las restricciones de secuencia, calculada por instancia (ver abajo)} \NR
\stopalign
\stopformula
\stopsubsection

\startsubsection[title={Parámetros derivados}]

Se calculan a partir de los datos, antes de resolver:
\startformula
\startalign[n=2, align={right,left}]
    \NC a_{bi} \NC = \startmathcases
                    \NC 1 \MC \text{si alguna orden del lote } b \text{ tiene productos en } i \text{, es decir, } \sum_{n \in \mathcal{N}_b} A_{ni} \ge 1 \NR
                    \NC 0 \MC \text{en otro caso} \NR
                 \stopmathcases \NR
    \NC L_b \NC = \sum_{i \in \mathcal{V}^{int}} a_{bi} \quad \text{(cantidad de locaciones que el lote } b \text{ debe visitar)} \NR
    \NC F_b \NC = T^{setup} + T^{pick} \cdot \sum_{n \in \mathcal{N}_b} C_n \quad \text{(tiempo fijo del lote: no depende de la ruta ni del picker)} \NR
\stopalign
\stopformula

El tiempo que tarda el picker $p$ en ejecutar el lote $b$ es su tiempo fijo más el tiempo de
ruta ajustado por su velocidad:
\startformula
T^{exec}_{pb} = F_b + \frac{r_b}{V_p}
\stopformula
donde $r_b$ es el tiempo de ruta del lote (traslado + servicio a velocidad de referencia). En
el modelo monolítico $r_b$ es una variable; en el de dos etapas es un dato.

$M^{seq}$ se calcula para cada instancia como una cota del tiempo total: lo que tardaría el
picker más lento en ejecutar todos los lotes seguidos, con una ruta de peor caso. Se redondea al
entero y tiene un piso de 500:
\startformula
\startalign[n=2, align={right,left}]
    \NC M^{seq} \NC = \max\left(500,\; \sum_{b \in \mathcal{B}} \left( F_b + \frac{R^{max}}{\min_{p} V_p} \right) + 1 \right) \NR
    \NC R^{max} \NC = T^{travel} \cdot |\mathcal{V}| \cdot \max_{(i,j) \in \mathcal{E}} D_{ij} \;+\; |\mathcal{V}| \cdot \max_{i \in \mathcal{V}} S_i \NR
\stopalign
\stopformula
\stopsubsection

\startsubsection[title={Cómo leer las restricciones con Big-M}]

Varias restricciones solo deben valer cuando una variable binaria vale 1, por ejemplo cuando
el lote $b$ está asignado al picker $p$ ($x_{pb} = 1$). Se escriben así:
\startformula
\text{expresión} \;\le\; \text{cota} + M \cdot (1 - x_{pb})
\stopformula
\startitemize
\item Si $x_{pb} = 1$, el término $M \cdot (1 - x_{pb})$ vale 0 y la restricción vale tal cual.
\item Si $x_{pb} = 0$, se suma $M$, que es tan grande que la restricción se cumple siempre: queda
desactivada.
\stopitemize
Cuando la condición involucra varias binarias, el término es $M \cdot (k - \text{suma de las
binarias})$: vale 0 solo si las $k$ binarias valen 1.
\stopsubsection

\stopsection


\page
%%%%%%%%%%%%%%%%%%%%%%%%%%%%%%%%%%%%%%%%%%%%%%%%%%%%%%%%%%%%%%%%%%%%
%%%%%%%%%%%%%%%%%%%%%%%%%%%%%%%%%%%%%%%%%%%%%%%%%%%%%%%%%%%%%%%%%%%%
\startsection[title={OPP v5.1 — Formulación monolítica}]

Un único MIP decide a la vez la ruta de cada lote, a qué picker se asigna y en qué orden lo
ejecuta.

\startsubsection[title={Variables}]

{\bf Ruteo}
\startformula
\startalign[n=2, align={right,left}]
    \NC z_{ijb} \in \{0,1\} \NC \quad 1 \text{ si la ruta del lote } b \text{ usa el arco } (i,j) \NR
    \NC q_{ijb} \ge 0 \NC \quad \text{cuántas locaciones del lote } b \text{ quedan por visitar al recorrer el arco } (i,j) \NR
    \NC r_b \ge 0 \NC \quad \text{tiempo de ruta del lote } b \text{ (traslado + servicio)} \NR
\stopalign
\stopformula

{\bf Asignación y secuencia}
\startformula
\startalign[n=2, align={right,left}]
    \NC x_{pb} \in \{0,1\} \NC \quad 1 \text{ si el lote } b \text{ se asigna al picker } p \NR
    \NC y_{bk} \in \{0,1\} \NC \quad \text{para } b < k \text{: } 1 \text{ si } b \text{ se ejecuta antes que } k \text{ (si comparten picker)} \NR
\stopalign
\stopformula

{\bf Tiempos}
\startformula
\startalign[n=2, align={right,left}]
    \NC u_{pb} \in [0,\; T^{cold} + \varepsilon] \NC \quad \text{tiempo de ejecución del lote } b \text{ si lo hace el picker } p \NR
    \NC c_{pb} \ge 0 \NC \quad \text{momento en que termina el lote } b \text{ en la agenda del picker } p \NR
    \NC c_b \ge 0 \NC \quad \text{momento en que termina el lote } b \NR
    \NC ta_n \ge 0,\; ea_n \ge 0 \NC \quad \text{tardanza y adelanto de la orden } n \text{ respecto de } T^{due}_n \NR
\stopalign
\stopformula

La cota superior de $u_{pb}$ es la cadena de frío (ver (15)).
\stopsubsection

\startsubsection[title={Función objetivo}]

Minimizar la suma ponderada de los tiempos de completitud de los lotes y de las penalizaciones
por adelanto y tardanza de las órdenes:
\startformula
\min \;\; W^{c} \cdot \sum_{b \in \mathcal{B}} c_b \;+\; W^{pen} \cdot \left( W^{e} \cdot \sum_{n \in \mathcal{N}} ea_n + W^{t} \cdot \sum_{n \in \mathcal{N}} ta_n \right)
\stopformula

El tiempo de ruta no aparece en el objetivo: solo influye a través de los tiempos de
completitud y de la cadena de frío.
\stopsubsection

\startsubsection[title={Restricciones de ruteo}]

Cada lote $b$ recorre un camino simple desde $v_{start}$ hasta $v_{end}$ que pasa por todas sus
locaciones requeridas ($a_{bi} = 1$). Puede atravesar otras locaciones, pero ninguna dos veces.

\startsubsubsection[title={(1)--(4) Inicio y fin de la ruta}]
La ruta sale una vez del depósito y llega una vez al nodo final. Si el depósito y el nodo final
son distintos, además no se vuelve al depósito ni se sale del nodo final.
\startformula
\startalign[n=3, align={right,left,left}]
    \NC (1) \NC \sum_{(i,j) \in \delta^{+}(v_{start})} z_{ijb} = 1 \NC \qquad \forall b \in \mathcal{B} \NR
    \NC (2) \NC \sum_{(i,j) \in \delta^{-}(v_{end})} z_{ijb} = 1 \NC \qquad \forall b \in \mathcal{B} \NR
    \NC (3) \NC \sum_{(i,j) \in \delta^{-}(v_{start})} z_{ijb} = 0 \NC \qquad \forall b \in \mathcal{B} \text{, si } v_{start} \ne v_{end} \NR
    \NC (4) \NC \sum_{(i,j) \in \delta^{+}(v_{end})} z_{ijb} = 0 \NC \qquad \forall b \in \mathcal{B} \text{, si } v_{start} \ne v_{end} \NR
\stopalign
\stopformula
\stopsubsubsection

\startsubsubsection[title={(5)--(9) Visitas y continuidad}]
(5) y (6): a cada locación requerida se entra y se sale. (7): en las locaciones intermedias, lo
que entra sale. (8) y (9): a cada locación se entra y se sale como máximo una vez, así que
ninguna se repite.
\startformula
\startalign[n=3, align={right,left,left}]
    \NC (5) \NC \sum_{(i,j) \in \delta^{-}(j)} z_{ijb} \ge a_{bj} \NC \qquad \forall b \in \mathcal{B},\; j \in \mathcal{V}^{int} \NR
    \NC (6) \NC \sum_{(j,i) \in \delta^{+}(j)} z_{jib} \ge a_{bj} \NC \qquad \forall b \in \mathcal{B},\; j \in \mathcal{V}^{int} \NR
    \NC (7) \NC \sum_{(g,i) \in \delta^{-}(i)} z_{gib} = \sum_{(i,j) \in \delta^{+}(i)} z_{ijb} \NC \qquad \forall b \in \mathcal{B},\; i \in \mathcal{V}^{int} \NR
    \NC (8) \NC \sum_{(j,i) \in \delta^{-}(i)} z_{jib} \le 1 \NC \qquad \forall b \in \mathcal{B},\; i \in \mathcal{V} \setminus \{v_{start}\} \NR
    \NC (9) \NC \sum_{(i,j) \in \delta^{+}(i)} z_{ijb} \le 1 \NC \qquad \forall b \in \mathcal{B},\; i \in \mathcal{V} \setminus \{v_{end}\} \NR
\stopalign
\stopformula
\stopsubsubsection

\startsubsubsection[title={(10)--(12) Eliminación de ciclos sueltos (Gavish-Graves)}]
(1)--(9) todavía permiten soluciones con un camino del depósito al final más un ciclo aparte
que visita algunas locaciones requeridas sin conectarse con el depósito. Para evitarlo se
envía un flujo $q$: sale del depósito con $L_b$ unidades, una por cada locación requerida, y
cada locación requerida consume una. Como el flujo solo circula por arcos usados (10) y
nace solo en el depósito (11), una locación requerida únicamente puede recibir su unidad si
está conectada al depósito por la ruta.
\startformula
\startalign[n=3, align={right,left,left}]
    \NC (10) \NC q_{ijb} \le L_b \cdot z_{ijb} \NC \qquad \forall b \in \mathcal{B},\; (i,j) \in \mathcal{E} \NR
    \NC (11) \NC \sum_{(i,j) \in \delta^{+}(v_{start})} q_{ijb} = L_b \NC \qquad \forall b \in \mathcal{B} \NR
    \NC (12) \NC \sum_{(g,i) \in \delta^{-}(i)} q_{gib} - \sum_{(i,j) \in \delta^{+}(i)} q_{ijb} = a_{bi} \NC \qquad \forall b \in \mathcal{B},\; i \in \mathcal{V}^{int} \NR
\stopalign
\stopformula
\stopsubsubsection

\startsubsubsection[title={(13) Tiempo de ruta}]
El tiempo de ruta es la distancia recorrida por $T^{travel}$ más el tiempo de servicio de cada
locación intermedia por la que se pasa (se cuenta una locación cuando la ruta sale de ella).
Sea $\mathcal{E}^{r} = \{(i,j) \in \mathcal{E} : i \ne v_{end},\; j \ne v_{start}\}$:
\startformula
r_b = T^{travel} \cdot \sum_{(i,j) \in \mathcal{E}^{r}} D_{ij} \, z_{ijb}
\;+\; \sum_{i \in \mathcal{V}^{int}} S_i \sum_{(i,j) \in \delta^{+}(i)} z_{ijb}
\qquad \forall b \in \mathcal{B}
\stopformula
Es decir, no suman los arcos que entran al depósito o salen del nodo final. Si $v_{start} = v_{end}$, esto
deja fuera también el primer y el último tramo de la ruta.
\stopsubsubsection
\stopsubsection

\startsubsection[title={(14)--(18) Asignación y tiempo de ejecución}]
(14): cada lote se asigna a exactamente un picker. (15) y (16): si el picker $p$ ejecuta el lote
$b$, $u_{pb}$ es exactamente su tiempo de ejecución $F_b + r_b / V_p$; si no, quedan
desactivadas. (17): el lote no termina antes de lo que tarda en ejecutarse. (18): el momento
de terminación del lote es el de la agenda del picker que lo ejecuta.
\startformula
\startalign[n=3, align={right,left,left}]
    \NC (14) \NC \sum_{p \in \mathcal{P}} x_{pb} = 1 \NC \qquad \forall b \in \mathcal{B} \NR
    \NC (15) \NC u_{pb} \ge F_b + \frac{r_b}{V_p} - M \cdot (1 - x_{pb}) \NC \qquad \forall p \in \mathcal{P},\; b \in \mathcal{B} \NR
    \NC (16) \NC u_{pb} \le F_b + \frac{r_b}{V_p} + M \cdot (1 - x_{pb}) \NC \qquad \forall p \in \mathcal{P},\; b \in \mathcal{B} \NR
    \NC (17) \NC c_{pb} \ge u_{pb} \NC \qquad \forall p \in \mathcal{P},\; b \in \mathcal{B} \NR
    \NC (18) \NC c_b \ge c_{pb} - M \cdot (1 - x_{pb}) \NC \qquad \forall p \in \mathcal{P},\; b \in \mathcal{B} \NR
\stopalign
\stopformula

{\bf Cadena de frío.} No hace falta una restricción aparte: está en la cota de la variable,
$u_{pb} \le T^{cold} + \varepsilon$. Si $x_{pb} = 1$, (15) obliga a
$u_{pb} \ge F_b + r_b / V_p$, así que el picker $p$ solo puede ejecutar el lote $b$ si
$F_b + r_b / V_p \le T^{cold} + \varepsilon$. Si $x_{pb} = 0$, $u_{pb}$ puede valer 0 y la cota
no restringe nada.
\stopsubsection

\startsubsection[title={(19)--(20) Secuencia sin solapamiento}]
Un picker ejecuta un lote por vez. Para cada picker $p$ y cada par de lotes $b < k$:
\startformula
\startalign[n=3, align={right,left,left}]
    \NC (19) \NC c_{pk} \ge c_{pb} + F_k + \frac{r_k}{V_p} - M^{seq} \cdot (3 - x_{pb} - x_{pk} - y_{bk}) \NC \qquad \forall p \in \mathcal{P},\; b < k \NR
    \NC (20) \NC c_{pb} \ge c_{pk} + F_b + \frac{r_b}{V_p} - M^{seq} \cdot (2 - x_{pb} - x_{pk} + y_{bk}) \NC \qquad \forall p \in \mathcal{P},\; b < k \NR
\stopalign
\stopformula
\startitemize
\item Si $p$ ejecuta los dos lotes y $y_{bk} = 1$ ($b$ va primero), solo vale (19): $k$ termina
al menos su tiempo de ejecución después de que termina $b$.
\item Si $p$ ejecuta los dos lotes y $y_{bk} = 0$ ($k$ va primero), solo vale (20), con los
papeles invertidos.
\item Si $p$ no ejecuta alguno de los dos, las dos quedan desactivadas.
\stopitemize
\stopsubsection

\startsubsection[title={(21)--(22) Tardanza y adelanto}]
Cada orden termina cuando termina su lote. Su tardanza es lo que se pasa de la fecha de
entrega y su adelanto, lo que le falta para llegar a ella:
\startformula
\startalign[n=3, align={right,left,left}]
    \NC (21) \NC ta_n \ge c_{pb} - T^{due}_n - M \cdot (1 - x_{pb}) \NC \qquad \forall p \in \mathcal{P},\; b \in \mathcal{B},\; n \in \mathcal{N}_b \NR
    \NC (22) \NC ea_n \ge T^{due}_n - c_{pb} - M \cdot (1 - x_{pb}) \NC \qquad \forall p \in \mathcal{P},\; b \in \mathcal{B},\; n \in \mathcal{N}_b \NR
\stopalign
\stopformula
Como $ta_n, ea_n \ge 0$ y el objetivo los minimiza, en el óptimo
$ta_n = \max(0,\, c_b - T^{due}_n)$ y $ea_n = \max(0,\, T^{due}_n - c_b)$. En el código ZIMPL el
término se escribe $M \cdot (2 - x_{pb} - B_{nb})$, que es igual porque solo se generan las
restricciones con $B_{nb} = 1$.
\stopsubsection

\startsubsection[title={Espera voluntaria}]
(17) y (19)--(20) son desigualdades: piden que un lote termine {\em no antes} de lo que tarda,
pero permiten declarar un $c_{pb}$ mayor. Cuando el adelanto pesa en el objetivo, al modelo le
puede convenir declarar un lote terminado más tarde para acercarlo a la fecha de entrega. Es
una propiedad de la formulación, compartida por las dos versiones de v5, y es la razón por la
que algunas instancias no se comparan en objetivo con las heurísticas (sección
\in[sec:comparacion]).
\stopsubsection

\stopsection


\page
%%%%%%%%%%%%%%%%%%%%%%%%%%%%%%%%%%%%%%%%%%%%%%%%%%%%%%%%%%%%%%%%%%%%
%%%%%%%%%%%%%%%%%%%%%%%%%%%%%%%%%%%%%%%%%%%%%%%%%%%%%%%%%%%%%%%%%%%%
\startsection[title={OPP v5.2 — Formulación en dos etapas}]

El mismo problema se resuelve con dos MIP, uno después del otro:
\startitemize[n]
\item {\bf Etapa 1}: calcula la ruta más corta de cada lote.
\item {\bf Etapa 2}: con esos tiempos de ruta como datos, decide asignación, secuencia y
cadena de frío.
\stopitemize

Se puede separar porque la ruta de un lote no depende de qué picker lo ejecuta ni del orden
entre lotes, y una ruta más corta nunca empeora el objetivo: solo adelanta las terminaciones y
facilita cumplir la cadena de frío.

\startsubsection[title={Etapa 1: ruteo de cada lote}]

{\bf Variables:} $z_{ijb}$, $q_{ijb}$ y $r_b$, con el mismo significado que en el monolítico.

{\bf Objetivo:} minimizar el tiempo total de ruta.
\startformula
\min \sum_{b \in \mathcal{B}} r_b
\stopformula
Como los lotes no comparten ninguna restricción, minimizar la suma es lo mismo que minimizar
cada $r_b$ por separado, y en la práctica cada lote se resuelve como un MIP independiente.

{\bf Restricciones:} exactamente (1)--(13) del monolítico. No usa pickers, pesos del objetivo
ni tiempos de preparación o picking.

{\bf Resultado:} el tiempo de ruta mínimo de cada lote, $R_b = r_b^{*}$, que pasa a la
Etapa~2 como dato.
\stopsubsection

\startsubsection[title={Etapa 2: asignación, secuencia y cadena de frío}]

\startsubsubsection[title={Datos que vienen de la Etapa 1}]
Con $R_b$ fijo, el tiempo de ejecución de cada par (picker, lote) es un número conocido antes
de resolver, y la cadena de frío se puede chequear de antemano:
\startformula
\startalign[n=2, align={right,left}]
    \NC T^{exec}_{pb} \NC = F_b + \frac{R_b}{V_p} \NR
    \NC \mathcal{PB}^{elig} \NC = \left\{ (p,b) \in \mathcal{P} \times \mathcal{B} : T^{exec}_{pb} \le T^{cold} + \varepsilon \right\} \NR
\stopalign
\stopformula
$\mathcal{PB}^{elig}$ son los pares en que el picker puede ejecutar el lote sin violar la
cadena de frío. Las variables $x_{pb}$ y $c_{pb}$ existen solo para esos pares, así que la
cadena de frío se cumple por construcción y no necesita ninguna restricción.
\stopsubsubsection

\startsubsubsection[title={Variables}]
\startformula
\startalign[n=2, align={right,left}]
    \NC x_{pb} \in \{0,1\} \NC \quad 1 \text{ si el lote } b \text{ se asigna al picker } p \text{, para } (p,b) \in \mathcal{PB}^{elig} \NR
    \NC c_{pb} \ge 0 \NC \quad \text{momento en que termina el lote } b \text{ en la agenda de } p \text{, para } (p,b) \in \mathcal{PB}^{elig} \NR
    \NC c_b \ge 0 \NC \quad \text{momento en que termina el lote } b \NR
    \NC y_{bk} \in \{0,1\} \NC \quad \text{para } b < k \text{: } 1 \text{ si } b \text{ se ejecuta antes que } k \NR
    \NC ta_n \ge 0,\; ea_n \ge 0 \NC \quad \text{tardanza y adelanto de la orden } n \NR
\stopalign
\stopformula
No hace falta la variable $u_{pb}$ del monolítico, porque el tiempo de ejecución ya es un dato.
\stopsubsubsection

\startsubsubsection[title={Función objetivo}]
La misma del monolítico:
\startformula
\min \;\; W^{c} \cdot \sum_{b \in \mathcal{B}} c_b \;+\; W^{pen} \cdot \left( W^{e} \cdot \sum_{n \in \mathcal{N}} ea_n + W^{t} \cdot \sum_{n \in \mathcal{N}} ta_n \right)
\stopformula
\stopsubsubsection

\startsubsubsection[title={(1)--(3) Asignación y terminación}]
(1): cada lote se asigna a exactamente un picker elegible. (2): si el picker $p$ ejecuta el lote
$b$, el lote no termina antes de su tiempo de ejecución; si no lo ejecuta, el lado izquierdo es
0 y no restringe. (3): el momento de terminación del lote es el de la agenda del picker que lo
ejecuta.
\startformula
\startalign[n=3, align={right,left,left}]
    \NC (1) \NC \sum_{p : (p,b) \in \mathcal{PB}^{elig}} x_{pb} = 1 \NC \qquad \forall b \in \mathcal{B} \NR
    \NC (2) \NC c_{pb} \ge T^{exec}_{pb} \cdot x_{pb} \NC \qquad \forall (p,b) \in \mathcal{PB}^{elig} \NR
    \NC (3) \NC c_b \ge c_{pb} - M \cdot (1 - x_{pb}) \NC \qquad \forall (p,b) \in \mathcal{PB}^{elig} \NR
\stopalign
\stopformula
(2) reemplaza a (15)--(17) del monolítico: como $T^{exec}_{pb}$ es un número, el producto
$T^{exec}_{pb} \cdot x_{pb}$ es lineal y no hace falta Big-M.
\stopsubsubsection

\startsubsubsection[title={(4)--(5) Secuencia sin solapamiento}]
Iguales a (19)--(20) del monolítico, con $T^{exec}_{pb}$ como dato. Se generan para cada picker
$p$ y cada par de lotes $b < k$ elegibles para $p$:
\startformula
\startalign[n=3, align={right,left,left}]
    \NC (4) \NC c_{pk} \ge c_{pb} + T^{exec}_{pk} - M^{seq} \cdot (3 - x_{pb} - x_{pk} - y_{bk}) \NC \qquad (p,b),(p,k) \in \mathcal{PB}^{elig},\; b < k \NR
    \NC (5) \NC c_{pb} \ge c_{pk} + T^{exec}_{pb} - M^{seq} \cdot (2 - x_{pb} - x_{pk} + y_{bk}) \NC \qquad (p,b),(p,k) \in \mathcal{PB}^{elig},\; b < k \NR
\stopalign
\stopformula
Se leen igual que en el monolítico: solo una de las dos está activa, según qué lote va primero,
y ambas se desactivan si $p$ no ejecuta los dos lotes.
\stopsubsubsection

\startsubsubsection[title={(6)--(7) Tardanza y adelanto}]
Iguales a (21)--(22) del monolítico, para los pares elegibles:
\startformula
\startalign[n=3, align={right,left,left}]
    \NC (6) \NC ta_n \ge c_{pb} - T^{due}_n - M \cdot (1 - x_{pb}) \NC \qquad \forall (p,b) \in \mathcal{PB}^{elig},\; n \in \mathcal{N}_b \NR
    \NC (7) \NC ea_n \ge T^{due}_n - c_{pb} - M \cdot (1 - x_{pb}) \NC \qquad \forall (p,b) \in \mathcal{PB}^{elig},\; n \in \mathcal{N}_b \NR
\stopalign
\stopformula
\stopsubsubsection

\startsubsubsection[title={(8)--(9) Dominio de las variables}]
\startformula
\startalign[n=3, align={right,left,left}]
    \NC (8) \NC ta_n,\; ea_n,\; c_b,\; c_{pb} \ge 0 \NC \qquad \forall n \in \mathcal{N},\; b \in \mathcal{B},\; (p,b) \in \mathcal{PB}^{elig} \NR
    \NC (9) \NC x_{pb},\; y_{bk} \in \{0,1\} \NC \qquad \forall (p,b) \in \mathcal{PB}^{elig},\; b < k \NR
\stopalign
\stopformula
\stopsubsubsection
\stopsubsection

\stopsection


\page
%%%%%%%%%%%%%%%%%%%%%%%%%%%%%%%%%%%%%%%%%%%%%%%%%%%%%%%%%%%%%%%%%%%%
\startsection[title={Diferencias entre las formulaciones}]

\bTABLE[split=yes, align=middle, frame=on, width=broad]
\setupTABLE[row][1][background=color, backgroundcolor=darkgray, style=bold]
\bTR \bTH Aspecto \eTH \bTH v5.1 Monolítico \eTH \bTH v5.2 Dos etapas \eTH \eTR
\bTR \bTD Número de MIPs \eTD \bTD 1 \eTD \bTD 2 (secuenciales) \eTD \eTR
\bTR \bTD Tiempo de ruta ($r_b$ / $R_b$) \eTD \bTD Variable de decisión \eTD \bTD Dato (resultado de Etapa 1) \eTD \eTR
\bTR \bTD $T^{exec}_{pb}$ \eTD \bTD Expresión con variable ($F_b + r_b / V_p$) \eTD \bTD Parámetro precalculado \eTD \eTR
\bTR \bTD Variable $u_{pb}$ \eTD \bTD Sí (con Big-M en (15)--(16)) \eTD \bTD No (innecesaria) \eTD \eTR
\bTR \bTD Cadena de frío \eTD \bTD Cota superior: $u_{pb} \le T^{cold} + \varepsilon$ \eTD \bTD Dominio: $\mathcal{PB}^{elig}$ \eTD \eTR
\bTR \bTD Dominio de $x_{pb}$ \eTD \bTD $\mathcal{P} \times \mathcal{B}$ (todos los pares) \eTD \bTD $\mathcal{PB}^{elig}$ (solo elegibles) \eTD \eTR
\bTR \bTD Piso de completitud \eTD \bTD $u_{pb} \le c_{pb}$ (17) \eTD \bTD $T^{exec}_{pb} \cdot x_{pb} \le c_{pb}$ (2) \eTD \eTR
\bTR \bTD Ruta garantizada mínima \eTD \bTD No (puede elegir ruta no mínima) \eTD \bTD Sí (Etapa 1 minimiza $r_b$) \eTD \eTR
\bTR \bTD Tamaño del MIP principal \eTD \bTD Grande (ruteo + scheduling juntos) \eTD \bTD Menor (cada etapa por separado) \eTD \eTR
\eTABLE

\stopsection


\page
%%%%%%%%%%%%%%%%%%%%%%%%%%%%%%%%%%%%%%%%%%%%%%%%%%%%%%%%%%%%%%%%%%%%
%%%%%%%%%%%%%%%%%%%%%%%%%%%%%%%%%%%%%%%%%%%%%%%%%%%%%%%%%%%%%%%%%%%%
\startsection[title={Heurísticas: HC-4 y ACO-4}, reference=sec:heuristicas]

OPP\_hc\_4 y OPP\_aco\_4 resuelven el mismo problema que v5, con la misma notación de la
sección anterior: ruteo de cada lote, asignación de lotes a pickers y secuencia de cada
picker, con velocidades heterogéneas y cadena de frío. No agregan congestión entre pickers
(a diferencia de OPP\_hc\_3 y OPP\_aco\_3), porque v5 tampoco la modela.

\startsubsection[title={Representación común de una solución}]

Una solución $S$ de cualquiera de las dos heurísticas tiene tres componentes:
\startitemize
\item {\bf Ruta} $\rho_b$ de cada lote $b$: un camino simple de $v_{start}$ a $v_{end}$ por
arcos de $\mathcal{E}$ que visita todas las locaciones requeridas del lote. Como en v5, no se
repite ningún vértice (restricciones (8)--(9) del monolítico).
\item {\bf Asignación} $p(b) \in \mathcal{P}$ de cada lote a un picker (equivale a $x_{pb}$).
\item {\bf Secuencia} $\sigma_p = (b_1, \dots, b_k)$ de los lotes asignados a cada picker
(equivale a $y_{bk}$).
\stopitemize

El tiempo de ruta y el tiempo de ejecución se calculan como en v5, pero sobre la ruta
concreta de la heurística:
\startformula
\startalign
    \NC R_b(\rho_b) \NC = T^{travel} \cdot \sum_{(i,j) \in \rho_b} D_{ij} \;+\; \sum_{i \in \rho_b \setminus \{v_{start},v_{end}\}} S_i \NR
    \NC T^{exec}_{pb} \NC = F_b + \frac{R_b(\rho_b)}{V_p} \NR
\stopalign
\stopformula

Los lotes de cada picker se ejecutan uno detrás del otro, sin esperas: si
$\sigma_p = (b_1, \dots, b_k)$,
\startformula
c_{b_1} = T^{exec}_{p b_1}, \qquad c_{b_j} = c_{b_{j-1}} + T^{exec}_{p b_j} \quad (j = 2, \dots, k)
\stopformula
y para cada orden $n$ del lote $b$, $ta_n = \max(0,\, c_b - T^{due}_n)$ y
$ea_n = \max(0,\, T^{due}_n - c_b)$. El objetivo $f(S)$ es exactamente el de v5:
\startformula
f(S) = W^{c} \cdot \sum_{b \in \mathcal{B}} c_b + W^{pen} \cdot \left( W^{e} \cdot \sum_{n \in \mathcal{N}} ea_n + W^{t} \cdot \sum_{n \in \mathcal{N}} ta_n \right)
\stopformula

Esta es la única diferencia de modelado con v5: en v5 la restricción de completitud es una
desigualdad ($c_{pb} \ge T^{exec}_{pb}$ y las de secuenciación), así que el modelo puede
declarar un lote terminado más tarde de lo que termina para reducir el adelanto
({\em espera voluntaria}). Las heurísticas reportan el tiempo real de terminación.

\stopsubsection

\startsubsection[title={Cadena de frío: dura en la respuesta, blanda en la búsqueda}]

Un par $(p,b)$ es elegible si $T^{exec}_{pb} \le T^{cold} + \varepsilon$, como en
$\mathcal{PB}^{elig}$ de v5.2. La solución que reporta una heurística cumple siempre esa
condición para todos sus lotes. Durante la búsqueda, en cambio, se permite pasar por
soluciones que la violan, penalizadas por el exceso total:
\startformula
\startalign
    \NC \Delta^{cold}(S) \NC = \sum_{b \in \mathcal{B}} \max\left(0,\; T^{exec}_{p(b)\,b} - T^{cold} - \varepsilon\right) \NR
    \NC g(S) \NC = f(S) + 1000 \cdot \Delta^{cold}(S) \NR
\stopalign
\stopformula

$g$ decide qué movimientos se aceptan; la mejor solución se guarda aparte, solo entre las que
tienen $\Delta^{cold}(S) = 0$, y se mide con $f$. Así una ruta inicial demasiado larga para el
límite de frío no descarta la instancia: la búsqueda puede acortarla o pasar el lote a un
picker más rápido.

Antes de buscar, las dos heurísticas aplican un único chequeo de infactibilidad que es exacto,
no una estimación: con $R_b = 0$, el mejor caso posible, $T^{exec}_{pb} = F_b$ para todo
picker, así que si $F_b > T^{cold} + \varepsilon$ ninguna ruta ni ningún picker puede cumplir
el frío y la instancia se declara infactible. En cualquier otro caso
la búsqueda corre completa; si termina sin ninguna solución que cumpla el frío, reporta
infactible. Por eso una heurística {\em no prueba} infactibilidad: solo no encuentra solución.

\stopsubsection

\startsubsection[title={OPP\_hc\_4: Hill Climbing}]

Búsqueda local estocástica de primera mejora: en cada iteración se elige {\em un} movimiento y
{\em un} vecino al azar, y se acepta si baja $g$. No se recorre toda la vecindad antes de
moverse, salvo dentro del movimiento de ruta.

\starttyping[option=none]
ALGORITMO HC4(params)
    si algún lote no es elegible para ningún picker con R_b = 0:
        retornar infeasible                       # chequeo exacto

    S <- SolucionInicial()
    si S es None: retornar infeasible             # lote no ruteable
    mejor <- S si Delta_cold(S) = 0 si_no None
    iteracion <- 0 ; sin_mejora <- 0

    mientras iteracion < Max_it y sin_mejora < Max_imp
                                y tiempo < Max_time:
        iteracion <- iteracion + 1
        S' <- Vecino(S)
        si S' es None o g(S') >= g(S):
            sin_mejora <- sin_mejora + 1 ; continuar
        S <- S' ; sin_mejora <- 0
        si Delta_cold(S) = 0 y f(S) < f(mejor):
            mejor <- S

    retornar mejor (o infeasible si es None)


SUBRUTINA SolucionInicial()
    para cada lote b con órdenes:
        candidatas <- { vecino más cercano (Dijkstra),
                        inserción más lejana, hacia adelante y hacia atrás,
                        20 construcciones pick-to-pick al azar,
                          con peso (1/d)^2.5, alternando el sentido,
                          cada una mejorada con 2-opt }
        si no hay candidatas: candidatas <- { búsqueda best-first acotada }
        rho_b <- pulir(la candidata más corta)    # 2-opt, or-opt, re-ruteo
    asignar cada lote al picker elegible con menos lotes
        (si no tiene ninguno elegible: al que minimiza T_exec_pb)
    retornar S


SUBRUTINA Vecino(S)
    elegir al azar entre los movimientos aplicables:
    Ruta:        lote al azar; aplicar la mejor inversión 2-opt factible
                 de su ruta, medida por distancia + penalización de frío
    Reasignar:   lote al azar; pasarlo a otro picker al azar
    Intercambio: picker al azar con 2 o más lotes; intercambiar dos
                 posiciones de su secuencia
\stoptyping

Parámetros en la suite de stress: $Max\_it = 200$, $Max\_imp = 60$ y $Max\_time = 30$~s. El
tope de tiempo se mira entre iteraciones, así que no corta una construcción inicial en curso.

\stopsubsection

\startsubsection[title={OPP\_aco\_4: Ant Colony Optimization}]

Cada hormiga construye una solución completa: para cada lote, primero la ruta y después el
picker. Hay dos tablas de feromona: una de ruta, $\tau_{ij}$ entre pares de locaciones
requeridas, y una de asignación, $\tau^{P}_{bp}$ entre lotes y pickers.

{\bf Construcción de la ruta.} Es {\em pick-to-pick} (Chen et al., 2013): desde la última
locación visitada $i$, la hormiga elige la próxima locación requerida $r$ entre las que faltan
con probabilidad
\startformula
P(r) = \frac{\tau_{ir}^{\alpha} \cdot \eta_{ir}^{\beta}}{\sum_{s \text{ pendiente}} \tau_{is}^{\alpha} \cdot \eta_{is}^{\beta}}, \qquad \eta_{ir} = \frac{1}{d(i,r)}
\stopformula
y camina hasta ella por el camino más corto. Como v5 no permite repetir vértices, $d(i,r)$ y
el camino se recalculan en cada paso con Dijkstra, evitando los vértices ya usados. Cada lote
tiene hasta 10 intentos de construcción, alternando el sentido (desde $v_{start}$ o desde
$v_{end}$); si todos fallan se usa una ruta de respaldo por búsqueda best-first. Cada ruta
construida se mejora con 2-opt y re-ruteo de segmentos.

{\bf Elección del picker.} Por ruleta entre los pickers elegibles para la ruta recién
construida, con peso $(\tau^{P}_{bp})^{\alpha} \cdot (1 / (1 + \text{carga}_p))^{\beta}$, donde
la carga es la cantidad de lotes que la hormiga ya le asignó. Si ningún picker es elegible, la
ruleta usa todos. La secuencia de cada picker se ordena por la fecha de entrega más temprana
de sus órdenes.

{\bf Feromonas.} La feromona de ruta sigue el esquema MAX-MIN (Stützle y Hoos, 2000): empieza
en $\tau_{max} = 2 / \rho$, en cada iteración se evapora y solo depositan, con
$\Delta\tau = 1$, la ruta más corta de la iteración y la más corta hasta el momento de cada
lote; se mantiene en $[\tau_{max}/50,\; \tau_{max}]$. Aprende de la longitud de la ruta, no
del objetivo completo. La feromona de asignación se evapora y recibe $Q / g(S)$ de cada
hormiga, solo en los lotes que cumplen el frío, más un depósito doble de la mejor solución.

{\bf Búsqueda local sobre la mejor solución.} Dos agregados propios, no tomados de la
bibliografía, que se aplican cuando cambia la mejor solución:
\startitemize
\item {\em Recombinación}: reemplazar la ruta de cada lote por la más corta que haya
construido cualquier hormiga, manteniendo asignación y secuencia; se acepta si baja $f$.
\item {\em Búsqueda local de agenda}: con las rutas fijas, mover un lote a otro picker en
cualquier posición, intercambiar dos lotes de la misma secuencia o intercambiar dos lotes entre
pickers, hasta que ningún movimiento mejore. Son las vecindades de HC-4, recorridas en forma
exhaustiva.
\stopitemize

\starttyping[option=none]
ALGORITMO ACO4(params)
    si algún lote no es elegible para ningún picker con R_b = 0:
        retornar infeasible                       # chequeo exacto
    inicializar tau (MAX-MIN) y tau_P
    mejor <- None

    para iteracion en 1 .. Max_it, mientras tiempo < Max_time:
        para hormiga en 1 .. K:
            S <- ConstruirSolucion(tau, tau_P)
            registrar la ruta más corta de cada lote
        pulir las rutas que mejoraron su récord
        S* <- la hormiga de menor f entre las que cumplen el frío
        si f(S*) < f(mejor): mejor <- S*
        mejor <- Recombinar(mejor) ; mejor <- BusquedaLocalAgenda(mejor)
        actualizar tau (mejores rutas por lote) y tau_P (todas las hormigas)

    retornar mejor (o infeasible si es None)
\stoptyping

Parámetros en la suite de stress: $K = 15$ hormigas, $Max\_it = 60$, $\alpha = 1{,}0$,
$\beta = 2{,}5$, $\rho = 0{,}15$, $Q = 100$ y $Max\_time = 30$~min. Como en HC-4, el tope de
tiempo se mira entre hormigas e iteraciones; la construcción puede usar hasta la mitad y el
resto queda para pulir y para la búsqueda local.

\stopsubsection

\startsubsection[title={Los cuatro métodos lado a lado}]

\blank[medium]
\bTABLE[split=yes, align=middle, frame=on, width=broad]
\setupTABLE[row][1][background=color, backgroundcolor=darkgray, style=bold]
\bTR \bTH Aspecto \eTH \bTH v5.1 mono \eTH \bTH v5.2 split \eTH \bTH HC-4 \eTH \bTH ACO-4 \eTH \eTR
\bTR \bTD Tipo \eTD \bTD MIP (ZIMPL + SCIP) \eTD \bTD Dos MIP (ZIMPL + SCIP) \eTD \bTD Búsqueda local \eTD \bTD Colonia de hormigas + búsqueda local \eTD \eTR
\bTR \bTD Optimalidad \eTD \bTD Probada \eTD \bTD Probada \eTD \bTD No \eTD \bTD No \eTD \eTR
\bTR \bTD Ruta \eTD \bTD Variable, junto con la agenda \eTD \bTD Mínima por lote (Etapa 1) \eTD \bTD Construcción + 2-opt \eTD \bTD Hormigas + 2-opt + recombinación \eTD \eTR
\bTR \bTD Asignación y secuencia \eTD \bTD Variables $x_{pb}$, $y_{bk}$ \eTD \bTD Variables $x_{pb}$, $y_{bk}$ \eTD \bTD Movimientos aleatorios \eTD \bTD Ruleta + búsqueda local \eTD \eTR
\bTR \bTD Cadena de frío \eTD \bTD Cota de $u_{pb}$ \eTD \bTD Dominio $\mathcal{PB}^{elig}$ \eTD \bTD Dura al reportar, penalizada al buscar \eTD \bTD Dura al reportar, penalizada al buscar \eTD \eTR
\bTR \bTD Espera voluntaria \eTD \bTD Sí \eTD \bTD Sí \eTD \bTD No \eTD \bTD No \eTD \eTR
\bTR \bTD Infactibilidad \eTD \bTD Probada \eTD \bTD Probada \eTD \bTD No encuentra solución \eTD \bTD No encuentra solución \eTD \eTR
\bTR \bTD Determinismo \eTD \bTD Sí \eTD \bTD Sí \eTD \bTD Semilla \eTD \bTD Semilla \eTD \eTR
\eTABLE
\blank[medium]

\stopsubsection

\stopsection


\page
%%%%%%%%%%%%%%%%%%%%%%%%%%%%%%%%%%%%%%%%%%%%%%%%%%%%%%%%%%%%%%%%%%%%
%%%%%%%%%%%%%%%%%%%%%%%%%%%%%%%%%%%%%%%%%%%%%%%%%%%%%%%%%%%%%%%%%%%%
\startsection[title={Comparación de v5 con las heurísticas}, reference=sec:comparacion]

Esta sección resume los resultados. El detalle instancia por instancia está en
\hyphenatedurl{opp_v5_2_split_vs_heuristics_comparison.tex} (split contra las heurísticas) y en
\hyphenatedurl{opp_v5_1_mono_vs_opp_v5_2_split_comparison.tex} (mono contra split).

\startsubsection[title={Comparabilidad}]

Se comparan solo instancias en que los cuatro métodos resuelven exactamente los mismos datos.
En la suite funcional son las instancias \type{test_v5eq__*} de las heurísticas, que capturan
el escenario de cada test funcional de v5 en el momento en que v5 escribe sus \type{.dat}. En
stress son los 19 escenarios de la suite, que los cuatro métodos construyen con el mismo
generador y los mismos argumentos. Las reglas están en
\type{docs/specs/comparabilidad_v5_heuristicas.md} y la fuente de verdad es
\type{results/comparability/comparability_v5_vs_heuristics.csv}.

\blank[medium]
\bTABLE[split=yes, align=middle, frame=on, width=fit]
\setupTABLE[row][1][background=color, backgroundcolor=darkgray, style=bold]
\setupTABLE[column][4][width=6cm, align=flushleft]
\bTR \bTH Situación de v5 \eTH \bTH Estado \eTH \bTH Objetivo \eTH \bTH Motivo \eTH \eTR
\bTR \bTD Óptimo, sin espera \eTD \bTD Sí \eTD \bTD Sí \eTD \bTD Mismos datos, mismo objetivo \eTD \eTR
\bTR \bTD Óptimo, con espera voluntaria \eTD \bTD Sí \eTD \bTD No \eTD \bTD v5 obtiene menos objetivo por diseño del modelo, no por un plan mejor \eTD \eTR
\bTR \bTD Infactible \eTD \bTD Sí \eTD \bTD No \eTD \bTD Se verifica que la heurística tampoco encuentre solución \eTD \eTR
\bTR \bTD No termina (timeout) \eTD \bTD No \eTD \bTD No \eTD \bTD No hay referencia; se usa para escala y para el experimento a igual tiempo \eTD \eTR
\eTABLE
\blank[medium]

\blank[medium]
\bTABLE[split=yes, align=middle, frame=on, width=fit]
\setupTABLE[row][1][background=color, backgroundcolor=darkgray, style=bold]
\bTR \bTH  \eTH \bTH Funcional \eTH \bTH Stress \eTH \eTR
\bTR \bTD Instancias \eTD \bTD 70 \eTD \bTD 19 \eTD \eTR
\bTR \bTD Comparables en estado y objetivo \eTD \bTD 45 \eTD \bTD 9 \eTD \eTR
\bTR \bTD Comparables solo en estado \eTD \bTD 22 \eTD \bTD 0 \eTD \eTR
\bTR \bTD \quad v5 infactible \eTD \bTD 12 \eTD \bTD 0 \eTD \eTR
\bTR \bTD \quad v5 usa espera voluntaria \eTD \bTD 10 \eTD \bTD 0 \eTD \eTR
\bTR \bTD No comparables \eTD \bTD 3 \eTD \bTD 10 \eTD \eTR
\eTABLE
\blank[medium]

En stress no hay espera que descartar: en los escenarios 04 a 19 el peso por lote completado
($W^{c} = 27$) es mayor que las órdenes por lote, así que esperar nunca conviene, y en los
heredados 01 a 03 todas las órdenes llegan tarde. Los 10 no comparables son los que split no
termina en 40 minutos. El monolítico resuelve las mismas instancias que split y llega al mismo
óptimo en los 49 pares funcionales que resuelven los dos, pero en stress no termina ninguno de
los 19 en 40 minutos, así que la referencia de stress es split.

\stopsubsection

\startsubsection[title={Validez y estado}]

Cada solución de las heurísticas se verificó contra OPP\_v5\_1\_mono: se fijan sus valores en las
variables del monolítico y SCIP decide si son factibles, sin reescribir ninguna restricción
(\type{verification/run_compliance.py}). Ninguna de las 150 soluciones verificadas (56
funcionales y 19 de stress por heurística) viola una restricción de v5, y el objetivo
recalculado coincide con el reportado en todas.

En estado, las dos heurísticas coinciden con split en las 67 instancias funcionales
comparables en estado y en los 9 escenarios de stress comparables: no devuelven solución en
las 12 infactibles y siempre encuentran una en las factibles.

\stopsubsection

\startsubsection[title={Calidad y tiempo hasta terminar}]

Brecha $= 100 \cdot (z_{\rm heur} - z_{\rm v5}) / z_{\rm v5}$ sobre las instancias comparables
en objetivo, con cada método corriendo hasta su propio criterio de parada (split hasta probar el
óptimo, con tope de 40 minutos).

\blank[medium]
\bTABLE[split=yes, align=middle, frame=on, width=fit]
\setupTABLE[row][1][background=color, backgroundcolor=darkgray, style=bold]
\bTR \bTH Heurística \eTH \bTH Suite \eTH \bTH Instancias \eTH \bTH Óptimo \eTH \bTH Brecha media \eTH \bTH Máxima \eTH \bTH Split / heurística (mediana) \eTH \eTR
\bTR \bTD HC-4 \eTD \bTD Funcional \eTD \bTD 45 \eTD \bTD 40 \eTD \bTD 0,36 \% \eTD \bTD 5,88 \% \eTD \bTD 5,1× \eTD \eTR
\bTR \bTD ACO-4 \eTD \bTD Funcional \eTD \bTD 45 \eTD \bTD 45 \eTD \bTD 0,00 \% \eTD \bTD 0,00 \% \eTD \bTD 1,4× \eTD \eTR
\bTR \bTD HC-4 \eTD \bTD Stress \eTD \bTD 9 \eTD \bTD 1 \eTD \bTD 0,06 \% \eTD \bTD 0,21 \% \eTD \bTD 53,1× \eTD \eTR
\bTR \bTD ACO-4 \eTD \bTD Stress \eTD \bTD 9 \eTD \bTD 2 \eTD \bTD 0,01 \% \eTD \bTD 0,06 \% \eTD \bTD 1,7× \eTD \eTR
\eTABLE
\blank[medium]

En los escenarios de stress que split resuelve:

\blank[medium]
\bTABLE[split=yes, align=middle, frame=on, width=fit]
\setupTABLE[row][1][background=color, backgroundcolor=darkgray, style=bold]
\bTR \bTH Caso \eTH \bTH Lotes \eTH \bTH Grilla \eTH \bTH Split (s) \eTH \bTH HC-4 (s) \eTH \bTH Brecha \eTH \bTH ACO-4 (s) \eTH \bTH Brecha \eTH \eTR
\bTR \bTD 02 \eTD \bTD 10 \eTD \bTD 10×10 \eTD \bTD 33 \eTD \bTD 0,9 \eTD \bTD 0,00 \% \eTD \bTD 93,5 \eTD \bTD 0,00 \% \eTD \eTR
\bTR \bTD 04 \eTD \bTD 6 \eTD \bTD 7×7 \eTD \bTD 12 \eTD \bTD 0,5 \eTD \bTD 0,01 \% \eTD \bTD 12,9 \eTD \bTD 0,00 \% \eTD \eTR
\bTR \bTD 05 \eTD \bTD 6 \eTD \bTD 13×13 \eTD \bTD 108 \eTD \bTD 2,8 \eTD \bTD 0,05 \% \eTD \bTD 84,1 \eTD \bTD 0,01 \% \eTD \eTR
\bTR \bTD 06 \eTD \bTD 6 \eTD \bTD 14×14 \eTD \bTD 202 \eTD \bTD 3,8 \eTD \bTD 0,04 \% \eTD \bTD 116,1 \eTD \bTD 0,01 \% \eTD \eTR
\bTR \bTD 07 \eTD \bTD 6 \eTD \bTD 15×15 \eTD \bTD 904 \eTD \bTD 4,3 \eTD \bTD 0,04 \% \eTD \bTD 156,1 \eTD \bTD 0,01 \% \eTD \eTR
\bTR \bTD 09 \eTD \bTD 6 \eTD \bTD 8×21 \eTD \bTD 224 \eTD \bTD 3,8 \eTD \bTD 0,06 \% \eTD \bTD 204,1 \eTD \bTD 0,00 \% \eTD \eTR
\bTR \bTD 10 \eTD \bTD 6 \eTD \bTD 7×7 \eTD \bTD 86 \eTD \bTD 0,5 \eTD \bTD 0,02 \% \eTD \bTD 15,3 \eTD \bTD 0,00 \% \eTD \eTR
\bTR \bTD 11 \eTD \bTD 6 \eTD \bTD 13×13 \eTD \bTD 139 \eTD \bTD 2,9 \eTD \bTD 0,21 \% \eTD \bTD 83,6 \eTD \bTD 0,06 \% \eTD \eTR
\bTR \bTD 15 \eTD \bTD 6 \eTD \bTD 10×10 \eTD \bTD 244 \eTD \bTD 1,5 \eTD \bTD 0,07 \% \eTD \bTD 52,5 \eTD \bTD 0,00 \% \eTD \eTR
\eTABLE
\blank[medium]

En los otros 10 escenarios (8 lotes o más, la grilla de 16×16 y el 100×100) split no termina
en 40 minutos y las dos heurísticas dan una solución válida: HC-4 en menos de 5~s y ACO-4 en
minutos, salvo el 01 (30 lotes en 100×100), que les lleva 17 y 35 minutos.

Esta comparación responde cuánto se pierde por no resolver en forma exacta, pero no es justa
en tiempo: split corre hasta {\em probar} el óptimo, y buena parte de ese tiempo puede ir a
cerrar la cota dual después de haber encontrado ya la mejor solución.

\stopsubsection

\startsubsection[title={A igual tiempo}, reference=sec:igual_tiempo]

La forma estándar de comparar una metaheurística con un solver exacto es darles el mismo
tiempo. En este experimento (\type{analysis/run_equal_time_study.py}) cada modelo de v5 recibe,
en cada escenario de stress, exactamente el tiempo total que tardó la heurística en ese
escenario, como límite de reloj de SCIP, y se registra la mejor solución que tiene al vencer
ese límite:
\startitemize
\item {\bf Split}: la Etapa 1 reparte el tiempo en partes iguales entre los lotes que faltan
rutear, y la Etapa 2 usa lo que queda. Solo hay solución completa si todos los lotes tienen
ruta y la Etapa 2 encuentra una solución a tiempo. Una ruta cortada por tiempo no es
necesariamente la mínima.
\item {\bf Mono}: una sola resolución con todo el tiempo.
\stopitemize

El tiempo se mide desde el inicio del pipeline (escribir los \type{.dat}, compilar con ZIMPL
y resolver), igual que el de la heurística; como el límite de SCIP cubre solo su resolución, el
tiempo transcurrido puede pasar un poco el presupuesto. Los resultados están en
\type{results/equal_time/}. La columna es $100 \cdot (z_{\rm heur} - z_{\rm v5}) / z_{\rm v5}$:
negativa significa que la heurística es mejor.

{\bf Con el tiempo de HC-4} (entre 0,5 y 5,7~s, salvo 1.209~s en el 01), ni split ni mono
obtienen una solución en ninguno de los 19 escenarios: en split, la Etapa 1 no llega a rutear
todos los lotes; el monolítico no encuentra ninguna solución entera. En el tiempo que HC-4
usa para dar una solución a menos de 0,21~\% del óptimo, v5 no da ninguna.

{\bf Con el tiempo de ACO-4:}

\blank[medium]
\bTABLE[split=yes, align=middle, frame=on, width=fit]
\setupTABLE[row][1][background=color, backgroundcolor=darkgray, style=bold]
\bTR \bTH Caso \eTH \bTH Tiempo (s) \eTH \bTH ACO-4 \eTH \bTH Split \eTH \bTH ACO-4 contra split \eTH \bTH Mono \eTH \bTH ACO-4 contra mono \eTH \eTR
\bTR \bTD 01 \eTD \bTD 1.877,8 \eTD \bTD 917.328,0 \eTD \bTD sin solución \eTD \bTD -- \eTD \bTD sin solución \eTD \bTD -- \eTD \eTR
\bTR \bTD 02 \eTD \bTD 31,3 \eTD \bTD 29.034,0 \eTD \bTD 29.262,0 \eTD \bTD -0,78 \% \eTD \bTD sin solución \eTD \bTD -- \eTD \eTR
\bTR \bTD 03 \eTD \bTD 412,1 \eTD \bTD 1.752.076,0 \eTD \bTD 1.763.704,0 \eTD \bTD -0,66 \% \eTD \bTD sin solución \eTD \bTD -- \eTD \eTR
\bTR \bTD 04 \eTD \bTD 15,1 \eTD \bTD 149.221,0 \eTD \bTD 149.218,9 \eTD \bTD 0,001 \% \eTD \bTD 157.569,7 \eTD \bTD -5,30 \% \eTD \eTR
\bTR \bTD 05 \eTD \bTD 95,0 \eTD \bTD 241.273,5 \eTD \bTD 241.398,5 \eTD \bTD -0,05 \% \eTD \bTD sin solución \eTD \bTD -- \eTD \eTR
\bTR \bTD 06 \eTD \bTD 109,8 \eTD \bTD 256.642,2 \eTD \bTD 256.791,0 \eTD \bTD -0,06 \% \eTD \bTD sin solución \eTD \bTD -- \eTD \eTR
\bTR \bTD 07 \eTD \bTD 194,6 \eTD \bTD 271.914,3 \eTD \bTD 272.016,8 \eTD \bTD -0,04 \% \eTD \bTD sin solución \eTD \bTD -- \eTD \eTR
\bTR \bTD 08 \eTD \bTD 129,5 \eTD \bTD 287.199,8 \eTD \bTD sin solución \eTD \bTD -- \eTD \bTD sin solución \eTD \bTD -- \eTD \eTR
\bTR \bTD 09 \eTD \bTD 223,9 \eTD \bTD 275.046,8 \eTD \bTD sin solución \eTD \bTD -- \eTD \bTD sin solución \eTD \bTD -- \eTD \eTR
\bTR \bTD 10 \eTD \bTD 15,2 \eTD \bTD 112.082,8 \eTD \bTD 112.083,2 \eTD \bTD -0,0004 \% \eTD \bTD sin solución \eTD \bTD -- \eTD \eTR
\bTR \bTD 11 \eTD \bTD 94,1 \eTD \bTD 182.338,5 \eTD \bTD sin solución \eTD \bTD -- \eTD \bTD sin solución \eTD \bTD -- \eTD \eTR
\bTR \bTD 12 \eTD \bTD 29,8 \eTD \bTD 198.905,5 \eTD \bTD 198.918,8 \eTD \bTD -0,007 \% \eTD \bTD 222.183,9 \eTD \bTD -10,48 \% \eTD \eTR
\bTR \bTD 13 \eTD \bTD 42,0 \eTD \bTD 248.696,3 \eTD \bTD 248.715,4 \eTD \bTD -0,008 \% \eTD \bTD sin solución \eTD \bTD -- \eTD \eTR
\bTR \bTD 14 \eTD \bTD 43,8 \eTD \bTD 298.381,2 \eTD \bTD 426.194,9 \eTD \bTD -29,99 \% \eTD \bTD sin solución \eTD \bTD -- \eTD \eTR
\bTR \bTD 15 \eTD \bTD 52,9 \eTD \bTD 97.922,8 \eTD \bTD 98.018,5 \eTD \bTD -0,10 \% \eTD \bTD sin solución \eTD \bTD -- \eTD \eTR
\bTR \bTD 16 \eTD \bTD 97,1 \eTD \bTD 195.901,5 \eTD \bTD 393.929,1 \eTD \bTD -50,27 \% \eTD \bTD sin solución \eTD \bTD -- \eTD \eTR
\bTR \bTD 17 \eTD \bTD 137,1 \eTD \bTD 305.798,8 \eTD \bTD 711.264,6 \eTD \bTD -57,01 \% \eTD \bTD sin solución \eTD \bTD -- \eTD \eTR
\bTR \bTD 18 \eTD \bTD 188,2 \eTD \bTD 439.821,4 \eTD \bTD 1.149.998,9 \eTD \bTD -61,75 \% \eTD \bTD sin solución \eTD \bTD -- \eTD \eTR
\bTR \bTD 19 \eTD \bTD 365,5 \eTD \bTD 780.799,0 \eTD \bTD 1.879.853,9 \eTD \bTD -58,47 \% \eTD \bTD sin solución \eTD \bTD -- \eTD \eTR
\eTABLE
\blank[medium]

\startitemize
\item {\bf Split} da solución en 15 de los 19 escenarios. ACO-4 es igual o mejor en 14 de esos
15; la excepción es el 04, donde split queda 0,001~\% por debajo. En 8 la diferencia es de
0,1~\% o menos, en el 02 y el 03 es de menos de 1~\%, y en los 5 escenarios con 12 lotes o más
(14, 16 a 19) split queda entre 30~\% y 62~\% por encima: con el tiempo repartido entre muchos
lotes, la Etapa 1 corta rutas lejos de la mínima, y la Etapa 2 arrastra esas rutas.
\item Split no da solución en el 01, el 08 y el 09 (la Etapa 1 no rutea todos los lotes) ni en
el 11, donde las rutas cortadas por tiempo hacen que algún lote no tenga ningún picker elegible
por cadena de frío. En el 09 la Etapa 1 excedió el presupuesto (864~s contra 224~s), porque el
límite se aplica a cada resolución y no al proceso completo.
\item {\bf Mono} da solución solo en el 04 y el 12, entre 5~\% y 10~\% por encima de ACO-4.
\stopitemize

Los presupuestos son los tiempos de la corrida de ACO-4 del 27 de septiembre, que no es la
misma corrida que la tabla anterior: por eso el tiempo y el objetivo de ACO-4 difieren un poco
(por ejemplo, 15,1~s contra 12,9~s en el 04, y 917.328 contra 988.252 en el 01).

\stopsubsection

\startsubsection[title={Conclusiones}]

\startitemize[n]
\item {\bf Mismo problema.} Las cuatro formulaciones comparten datos, restricciones y objetivo;
la única diferencia de modelado es la espera voluntaria, que v5 permite y las heurísticas no.
Las instancias donde v5 la usa se comparan solo en estado.
\item {\bf Las heurísticas son válidas.} Todas sus soluciones cumplen las restricciones de v5,
verificadas contra el monolítico, y coinciden con split en el estado de todas las instancias
comparables.
\item {\bf Cerca del óptimo.} ACO-4 alcanza el óptimo en las 45 funcionales comparables y queda
a 0,06~\% como máximo en stress; HC-4 falla en 5 funcionales de grilla (hasta 5,88~\%) y queda
a 0,21~\% como máximo en stress.
\item {\bf A igual tiempo, las heurísticas dominan.} Con el tiempo de HC-4, v5 no obtiene
ninguna solución. Con el de ACO-4, split empata o pierde en todos los escenarios salvo uno (por
0,001~\%), pierde por 30~\% a 62~\% con 12 lotes o más, y en 4 escenarios no da solución; el
monolítico casi nunca la da.
\item {\bf El valor de v5 es la prueba de optimalidad.} Hasta 6 lotes en grillas de hasta
15×15, split prueba el óptimo en minutos y sirve de referencia para medir a las heurísticas.
Por encima de ese tamaño, o cuando el tiempo de respuesta importa, las heurísticas son la
opción: ACO-4 cuando importa la calidad y HC-4 cuando importan los segundos.
\stopitemize

\stopsubsection

\stopsection

\stoptext
